% !PRINTED 365
\weberpage{365}
%<<CONT
each complex, the total number of the complexes is $14.27 : 6 = 63$.

\gesperrt{5. There are 63 \gesperrt{Steiner} complexes.}

If from the pairs of the complex (1) we pick out any other pair of pairs instead of $x_1 y_1, x_2 y_2$, then from it, following the scheme of (2) and (3), we can each time form two new complexes. Since there are 15 such pairs of pairs, we obtain 30 new complexes of the type (2), (3), which are all derived in the same way from the complex (1), and which are all different from one another.

If we now take any one of the double tangents $z_1$ newly entering in (2) and (3), we obtain two new complexes by starting from the two pairs $x_1 z_1, y_1 z_1$, and since we can choose $z_1$ in 16 different ways, there result in this way 32 new complexes, with which the total number of all the complexes is exhausted. But the distribution of the double tangents among the complexes $x_1 z_1, y_1 z_1$ must still be investigated more closely.

Without restricting generality, since if necessary we can interchange $x_1$ with $y_1$, we can make the assumption that $z_1$ occurs in the complex (2), and there in the pair $z_1 z_2$.

Then we obtain the two complexes:
\weq{4}{x_1 z_1,\quad x_2 z_2,\quad \ldots}
\weq{5}{y_1 z_1,\quad y_2 z_2,\quad \ldots}
We now consider the three complexes:
\weq{4}{x_1 z_1,\quad x_2 z_2,\quad \ldots}
\weq{2}{x_1 x_2,\quad z_1 z_2,\quad \ldots}
\weq{4a}{x_1 z_2,\quad x_2 z_1,\quad \ldots,}
which by Theorem 4. must contain all the double tangents, hence among them also $x_3, y_3$. But these do not occur in (2), and likewise they cannot both occur in (4) or both in (4a), since otherwise $x_1, x_3, y_3$ would have to be azygetic, whereas [by (1)] they are syzygetic. Since if need be we may interchange $x_3$ and $y_3$ in the notation, we can assume that $x_3$ occurs in (4), and indeed in a pair $x_3 z_3$.

Then, since $z_1 z_2 z_3$ are azygetic, $z_3$ cannot occur in the complex (2) and must consequently be contained in (3).
